% =========================================================
% Section 2 — Benchmark Setting and MineNPC-Task (merged, compact)
% =========================================================
\section[Benchmark Setting and \textsc{MineNPC-Task}]{Benchmark Setting and \textsc{MineNPC-Task}}
\label{sec:setting}

We instantiate open-world co-play in \emph{Minecraft} using a Mineflayer client so that actions are constrained to public, in-game interfaces. Concretely, the agent observes chat, its own inventory and equipment, and nearby entities/blocks within currently loaded chunks (a practical proxy for line of sight). Actions are issued through high-level skills like \emph{navigate, mine, craft, place, interact} implemented atop Mineflayer APIs. A bounded-knowledge policy forbids privileged capabilities (e.g., \texttt{/give}, \texttt{/teleport}), global map/seed introspection, and scans beyond loaded chunks; runs that violate this policy are invalidated.

\textbf{Task source and structure.} The \textsc{MineNPC-Task} suite is derived from goals observed during expert co-play rather than from synthetic prompts. From these observations we specify a lightweight \emph{task template} that the evaluation framework instantiates at run time. Given a user goal, the framework compiles a short plan (typically 3–5 subtasks) and represents each subtask as a compact record with five fields: name, dependencies, required parameters, a single targeted clarifying question issued only if a required parameter is missing, and a success criterion. Completion is judged by simple, machine-checkable validators that consume only bounded, in-world evidence—inventory/equipment and position deltas, nearby entities/blocks within loaded chunks, and a short window of recent chat—and return a pass/fail with a brief rationale. (Prompts and template examples appear in the Appendix.)

\textbf{Coverage.} The \textsc{MineNPC-Task} suite covers the everyday goals we observed most frequently: \emph{resource acquisition and logistics} (gather–craft chains, inventory constraints); \emph{navigation and retrieval} (recalling named landmarks and fetching/returning items); \emph{tooling and crafting} (multi-step recipes, tool selection with durability checks); \emph{construction} (layout-constrained builds using existing facilities); \emph{combat and safety} (simple loadout preconditions); and \emph{continuity/mixed-initiative} cases where an underspecified slot is bound via a brief clarification. Section~\ref{sec:framework} details how the evaluation framework enforces this template end-to-end and produces validator-backed outcomes.